\begin{enigme}[Belle figure]

Construis une figure analogue à partir d'un triangle $ABC$ isocèle de sommet principal $A$ tel que $BC = 10$ cm et $AC = 14$ cm :

\begin{center} \includegraphics[width=3.2cm]{t_artistique} \end{center}

 \end{enigme}
 
 
%%%%%%%%%%%%%%%%%%%%%%%%%%%%%%%%%%%%%%%%%%%%%%%%%%%%%%%%%%%%%

\begin{enigme}[Artistes en géométrie]

\begin{enumerate}
 \item Recherche des informations sur le peintre Pietr Mondrian et notamment sur ses œuvres peintes à Paris.
 \item Quelles figures géométriques sont souvent visibles dans ses toiles ?
 \item À la manière de Mondrian, sur une feuille blanche, trace un cadre avec, à l'intérieur, des droites parallèles verticales et horizontales. Puis colorie en t'inspirant des œuvres de cet artiste.
 \item L'artiste Vassily Kandinsky a aussi travaillé à partir de figures géométriques. Cite le nom de certaines de ses œuvres.
 \item Recherche d'autres artistes ayant travaillé avec des figures géométriques.
 \end{enumerate}
 
 \end{enigme}


